% Appendix: proofs for Section 2 (complementarity model).  Requires \appendix before inclusion.
\section{Proofs for the complementarity model}
\label{app:eps}

Throughout, agent~1 has type $(x,1-x)$, agent~2 has $(y,1-y)$, and $u_v(z_1,z_2)=vz_1+(1-v)z_2+\varepsilon z_1z_2$.

\paragraph{P1 (first best).} Giving agent~1 the bundle $(z_1,z_2)$ yields total welfare
$1+(x-y)(z_1-z_2)+\varepsilon[z_1z_2+(1-z_1)(1-z_2)]$, which is affine in each coordinate separately
and hence maximised at a corner of $[0,1]^2$: both items to one agent gives $1+\varepsilon$, the
valuation-based split gives $1+|x-y|$, the reverse split less. An average of allocations never
exceeds the best of them, so randomisation does not help. \qed

\paragraph{P2 (fixed menus).} The dictator's report only selects an element of the fixed menu
$M$; any misreport replaces her favourite element of $M$ by another element of $M$. The other
agent's report does not enter the dictator branch at all, and in the other branch the roles are
reversed. Hence truth-telling is dominant for both agents, for any utility function. \qed

\paragraph{P3 (dictator's problem).} Along the frontier $t\mapsto(Q(t),Q(1-t))$ the objective is
$F_v(t)=vQ(t)+(1-v)Q(1-t)+\varepsilon P(t)$ with $P=Q(t)Q(1-t)$. $Q''<0$, and
$P''=Q''(t)Q(1-t)-2Q'(t)Q'(1-t)+Q(t)Q''(1-t)<0$, so $F_v$ is strictly concave and the optimal spend
$t^\star(v)$ is unique. $\partial^2F/\partial t\,\partial v=Q'(t)+Q'(1-t)>0$, so by Topkis' theorem
$t^\star(v)$ is non-decreasing: the share of item~1 is non-decreasing in the type and that of
item~2 non-increasing. For nested lotteries the objective is $vm_1+(1-v)\varphi(m_1)+\varepsilon\min(m_1,\varphi(m_1))$
along a concave frontier $\varphi$; a minimum of concave functions is concave and the increasing
differences persist, so the same conclusions hold. \qed

\paragraph{P4 (certified brackets).} On a cell $[x_0,x_0+h]\times[y_0,y_0+h]$ the branch
``agent~1 dictates'' is bounded below by (i) $u_x(q(x))\ge u_x(q(x_0))$, since $q(x_0)$ is on the
menu, and the right side is affine in $x$ so it is at least its minimum over the two $x$-endpoints;
(ii) $u_y(1-q(x))\ge y_0(1-q_1(x_0+h))+(1-y_0-h)(1-q_2(x_0))+\varepsilon(1-q_1(x_0+h))(1-q_2(x_0))$, by the
monotonicity of $q$ and non-negativity of the coefficients (for nested lotteries replace the
product by $1-\max(q_1(x_0+h),q_2(x_0))$). Symmetrically for the other branch; the denominator is
at most $1+\max(\varepsilon,x_0+h-y_0)$. Exact ratios at cell corners bound the minimum from above,
and the adaptive refinement of Section~1 applies unchanged. \qed

\paragraph{P5 (strategic play against the naive rule).} Under the naive rule a dictator who
reports $r$ receives the frontier point at spend $t_0(r)=\operatorname{clip}((r(a+2b)-b)/a,0,1)$;
$r\mapsto t_0(r)$ is onto $[0,1]$, and her report is irrelevant in the other branch. Her dominant
strategy is therefore any $r$ with $t_0(r)=t^\star(x)$, and the realised allocation is that of the
$\varepsilon$-aware mechanism. For the misreport bound, write $F^\varepsilon_x(t)=F^0_x(t)+\varepsilon P(t)$;
$t_0(x)$ maximises $F^0_x$, so for any $s$, $F^\varepsilon_x(s)-F^\varepsilon_x(t_0(x))\le\varepsilon[P(s)-P(t_0(x))]\le\varepsilon\max P=\varepsilon Q(\tfrac12)^2$,
and the coin halves it. \qed

\paragraph{P6 (nested lotteries).} With one uniform draw $U$ and ``item $j$ iff $U<m_j$'', the
marginals are $m_1,m_2$, $\mathbb E[z_1z_2]=\min(m_1,m_2)\ge m_1m_2$, and for the other agent
$\mathbb E[(1-z_1)(1-z_2)]=1-\max(m_1,m_2)\ge(1-m_1)(1-m_2)$. Both agents' $\varepsilon$-terms weakly
increase and the menu is still fixed, so P2 applies. As $\varepsilon\to\infty$ the dictator's
objective is dominated by $\varepsilon\min(m_1,m_2)$, maximised at the balanced point
$m_1=m_2=Q(\tfrac12)$: she receives everything with probability $Q(\tfrac12)$ and the other agent
otherwise, welfare $1+\varepsilon=\mathrm{OPT}$; numerically this already holds for every
$\varepsilon\ge1$. \qed

\paragraph{P7 (moment LP).} For a lottery over bundles, agent~1's expected utility is
$x\,m_1+(1-x)\,m_2+\varepsilon\,m_{12}$ and agent~2's is
$y(1-m_1)+(1-y)(1-m_2)+\varepsilon(1-m_1-m_2+m_{12})$, so only $(m_1,m_2,m_{12})$ matter. The set of
attainable triples is the convex hull of $\{(z_1,z_2,z_1z_2):z\in[0,1]^2\}$, which is the McCormick
envelope $m_{12}\ge0$, $m_{12}\ge m_1+m_2-1$, $m_{12}\le m_1$, $m_{12}\le m_2$ (the four inequalities are
valid, and every point satisfying them is a convex combination of the corners of the square). On
a type grid the best strategy-proof mechanism is thus an LP in these variables with linear
incentive, competitiveness ($W=1+(x-y)(m_1-m_2)+\varepsilon(1-m_1-m_2+2m_{12})\ge\alpha\,\mathrm{OPT}$) and
McCormick constraints; restricting types only relaxes constraints, so its value bounds every
strategy-proof mechanism on the continuum. Utilities are affine in the own type, so the
adjacent-type reduction of Section~1 applies. \qed

\paragraph{Theorem B.} Let the dictator of type $x$ choose the lottery with moments
$m=(m_1,m_2,m_{12})$, $m_{12}\le\min(m_1,m_2)$, and let the other agent have type $y$. The branch
welfare is
\[
B(x,y,m)=(1+\varepsilon)+(x-y-\varepsilon)m_1+(y-x-\varepsilon)m_2+2\varepsilon m_{12}.
\]
\emph{(i) Extreme profiles.} Write $m(1)=(a_1,a_2,a_{12})$, $m(0)=(c_1,c_2,c_{12})$ and
$p=a_1-a_2$, $q=c_2-c_1$. Using $a_{12}\le a_2$, $c_{12}\le c_1$,
\[
W(1,1)\le(1+\varepsilon)-\varepsilon|p|,\quad W(0,0)\le(1+\varepsilon)-\varepsilon|q|,\quad
W(1,0)=\tfrac12[B(1,0,m(1))+B(0,1,m(0))]\le(1+\varepsilon)+\tfrac12(1-\varepsilon)(p+q),
\]
with first bests $1+\varepsilon$, $1+\varepsilon$, $2$. The minimum of the three ratios is concave and
symmetric in $(p,q)$, hence maximised at $p=q=s$, where $(1+\varepsilon-\varepsilon s)/(1+\varepsilon)$ and
$[(1+\varepsilon)+(1-\varepsilon)s]/2$ are equalised at $s^\star=(1-\varepsilon^2)/(1+2\varepsilon-\varepsilon^2)$,
giving $\rho_B$. \emph{(ii) Middle type.} On the four profiles $(\tfrac12,0),(1,\tfrac12),(0,\tfrac12),(\tfrac12,1)$
(first best $\tfrac32$ when $\varepsilon\le\tfrac12$) the middle type's two branches sum to
$2(1+\varepsilon)-2\varepsilon(r_1+r_2)+4\varepsilon r_{12}\le2(1+\varepsilon)$, and each extreme type's cross
branch is at most $(1+\varepsilon)+\max(0,\tfrac12-\varepsilon)$; summing, $\sum W\le4(1+\varepsilon)+2(\tfrac12-\varepsilon)$,
so some profile has ratio at most $(5+2\varepsilon)/6=\rho_A$. \emph{Tightness.} For $\varepsilon\ge\tfrac12$
the menu $\{(s^\star,0),(0,s^\star)\}$ has, on same-side profiles, $|x-y|\le\tfrac12\le\varepsilon$ so
$\mathrm{OPT}=1+\varepsilon$ and ratio $(1+\varepsilon-\varepsilon s^\star)/(1+\varepsilon)=\rho_B$; on opposite sides
the ratio $[(1+\varepsilon)+(x-y-\varepsilon)s^\star]/(1+\max(\varepsilon,x-y))$ is decreasing in $x-y$ and equals
$\rho_B$ at $x-y=1$. Exact evaluation of this finite menu on the continuum confirms the value
($6/7$ at $\varepsilon=\tfrac12$, $28/31$ at $\tfrac34$). \qed
